% Fichier complété par tools/complete_missing_corrections.py.
% Source : SYS/SYS-01/SYS-01_ChaineInfo/538_Codeur/538_Codeur.tex
% Statut : rédigé (5 question(s)).
\begin{corrigebox}[Corrigé rédigé]
\CorrigeQuestion{1}
Un codeur incrémental optique comporte un disque à fentes, une source et un photodétecteur. La rotation produit des impulsions ; deux voies A et B en quadrature donnent le sens et une voie index fournit une référence par tour.

\CorrigeQuestion{2}
Avec 25 fentes et une seule voie, la résolution vaut 14,4 degrés par impulsion. En comptant les deux fronts, elle devient 7,2 degrés, puis 3,6 degrés avec deux voies et quatre fronts.

\CorrigeQuestion{3}
La fréquence des impulsions est égale au nombre des impulsions par tour multiplié par la vitesse en tours par minute, puis divisé par 60. La vitesse est donc égale à 60 fois la fréquence divisée par le nombre des impulsions par tour. Le sens se déduit de la succession des fronts A et B.

\CorrigeQuestion{4}
La position relative est égale au nombre algébrique des impulsions multiplié par la résolution. Une prise de référence sur la voie index est nécessaire après la mise sous tension.

\CorrigeQuestion{5}
La valeur minimale de quantification vaut un demi-pas, soit 7,2 degrés avec une seule voie et 25 fentes, avant prise en compte des erreurs optiques et mécaniques.
\end{corrigebox}
